% 结论
\chapter*{结论}
\addcontentsline{toc}{chapter}{结论}

本文从居家瑜伽练习的动作层面评估问题出发，实现了一套基于多模态大模型的动作评估辅助分析系统，其具体表现流程从一个用户上传视频开始，到最后形成一份可供查看保存的评估报告，中间经历了用户视频上传、姿态关键点抽取、几何规则分析、反馈多模态生成、结果展现，整个流程不是一次性达成的，而是经历了Python原型核验、Rust后端试做、Flask+原生JavaScript工程落地等几个阶段，笔者的想法也随之逐渐清晰，本课题真正的难点不在于算法而在如何将算法、接口、数据库、交互真正连成稳定的闭环。

具体实现方面，本文实现了注册登录、上传视频、后台执行异步任务调度、姿态特征提取、规则评分、多模态反馈、结果呈现、历史记录留存等几个主要功能，并且对上文中联调发现的真实问题进行了修复，包括角色字段映射冲突、上传阻塞、SQLite锁冲突、触发器列引用错误、空对象异常、代理变量干扰模型请求等。整个系统经过几轮测试，已经能完成本地场景下业务的基本流程，并且可以在部分异常场景下保持降级输出。

从论文写作和项目实现的关系看，本文并没有把系统描述为一个已经成熟的商业化训练平台，而是把重点放在“如何把多个技术环节组织成可运行闭环”上。姿态估计、几何规则和多模态反馈分别解决了不同层面的问题：姿态估计负责把视频转成可计算的关键点；几何规则负责把关键点转成可解释的偏差和得分；多模态模型负责把偏硬的数据结果转成用户更容易理解的建议文字。三个环节相互配合后，系统才真正具备从输入到反馈的完整意义。若只完成其中某一环，论文论证都会显得不够充分。

在开发过程中，笔者也逐渐认识到，工程系统的难点常常不在单个函数本身，而在不同模块之间的边界是否清楚。前端字段和后端字段不一致、数据库触发器引用不严谨、模型接口受代理变量影响、任务状态没有及时回写，这些问题看似分散，但本质上都与接口契约和状态管理有关。经过多轮修复后，系统在结构上比早期原型更清晰，也更容易说明每个模块承担的职责。这一点对于毕业设计答辩同样重要，因为它能够展示项目不是简单堆叠代码，而是经历了设计、验证和调整的过程。

本文工作的另一个收获，是对“可解释反馈”的理解更加具体。对于运动训练类系统而言，用户并不只需要一个分数，也不只需要一句笼统评价，而是需要知道偏差来自哪里、为什么会影响动作质量、下一次应该怎样调整。规则评估提供了偏差定位，多模态模型提供了自然语言表达，前端展示则把结果组织成用户可阅读的形式。三者结合以后，系统的反馈才不只是技术输出，而是具有一定训练指导意义的信息。

同时，本文也保留了对系统局限性的说明。当前系统样本规模有限，动作标准库覆盖范围有限，单摄像头输入也难以完全解决三维姿态遮挡和角度恢复问题。因此，本文的结论主要限定在本地演示环境和小规模样本验证范围内。这样的边界并不削弱课题价值，反而能更准确地说明本文工作的性质：它完成的是一个可运行、可追踪、可扩展的原型系统，而不是面向大规模商业应用的最终产品。

如果从学习过程本身来看，本课题也使笔者对“系统实现”和“论文论证”之间的关系有了更直接的体会。系统开发不能只追求页面可见效果，论文写作也不能只停留在抽象描述。只有当开发中的真实问题被整理成论文章节中的设计依据，论文中的结构分析又能反过来解释系统实现路径时，毕业设计才形成较完整的闭环。本文中关于异步调度、数据库锁冲突、字段映射和模型降级的讨论，都是在真实开发问题基础上整理出来的。

同时，本课题也体现了轻量级工程方案的现实价值。相比直接构建复杂平台，本文选择 Flask、SQLite、原生 JavaScript 和本地文件系统，是为了在有限时间内优先完成可验证链路。这样的选择并不意味着技术含量较低，相反，它要求开发者更清楚地处理模块边界、异常流程和数据一致性。轻量系统如果结构混乱，同样会很快失控；而当它具备清晰的职责划分和排错路径时，也能较好地支撑毕业设计阶段的研究目标。

总体而言，本文完成的不是一个孤立算法实验，而是一套围绕居家瑜伽动作评估场景展开的综合系统实践。它把姿态估计、规则计算、多模态反馈、Web 服务、数据库存储和前端交互串联在一起，并通过测试说明该链路能够在本地环境中运行。虽然系统还有许多不足，但已经能够说明将多模态大模型用于运动训练反馈场景具有一定可行性，也为后续继续扩展体式库、优化可视化和提升稳定性提供了基础。

从毕业设计的出发点和目标看，本文已经完成了“做出一套可运行系统并说明其实现与验证过程”的主要任务，与此同时，系统还有很多后续改进的空间。动作标准库不够丰富，更大并发下性能边界还未系统地测量过，模型反馈稳定性还受依赖的外部服务的影响，后续如果继续的话，可从体式覆盖范围、任务调度、结果可视化、外部网络服务依赖这几个角度来继续完善。

首先，在体式覆盖范围方面，当前系统主要围绕有限的瑜伽动作进行验证，规则库仍然偏小。后续可以逐步扩展到更多常见体式，并为每个体式补充关键角度、容差范围、常见错误和纠正建议。更进一步，还可以根据不同训练阶段设置不同阈值，例如初学者允许更大的偏差范围，而进阶练习者则采用更严格的评估标准。这样可以让系统从“统一规则评估”逐渐走向“分层训练反馈”。

其次，在数据和评估方法方面，未来可以引入更多真实样本，建立不同拍摄角度、不同光照、不同身材比例和不同动作熟练度下的测试集合。当前系统更多是工程验证，能够说明链路可行，但还不能形成充分的统计结论。如果后续能够积累一定规模的标注样本，就可以对规则阈值、稳定性指标和模型反馈一致性做更严谨的评估，也可以比较不同姿态估计算法在瑜伽场景下的表现差异。

再次，在系统架构方面，当前后台任务主要依赖线程和轮询机制，足以支撑本地演示和小规模使用。若未来面向更多用户，就需要引入任务队列、持久化任务状态和更细粒度的并发控制，甚至可以把视频处理服务和 Web 接口服务拆分部署。数据库也可以从 SQLite 迁移到 MySQL 或 PostgreSQL，以获得更好的并发写入能力和运维工具支持。这样的演进并不会改变本文提出的基本流程，但会使系统更接近真实应用环境。

最后，在反馈表达和可视化方面，后续可以进一步增强结果呈现。除了文字建议，系统还可以在关键帧上标出偏差关节、用折线图展示角度随时间变化、用热力图说明身体左右侧差异，甚至生成简化的骨架动画帮助用户理解动作变化。相比单纯给出评分，这类可视化结果更容易帮助用户建立身体感知，也更符合运动训练类系统的使用习惯。与此同时，多模态模型反馈仍应保持受控，核心判断应继续建立在可追溯的规则和特征之上，避免让自然语言生成替代系统本身的可解释依据。

除此之外，后续还可以考虑增加更完整的用户反馈机制。当前系统更多是单向给出评估结果，用户是否采纳建议、下一次动作是否因此改善，还没有形成持续学习闭环。如果未来加入用户自评、教练复核或多次训练趋势分析，系统就可以从单次评估工具逐渐发展为训练记录与改进建议平台。届时，历史记录不只是存档，而可以用于观察用户动作稳定性是否提升、常见错误是否减少、不同体式的训练进度是否均衡。

在隐私和安全方面，后续也需要更加细致的设计。瑜伽训练视频涉及人体图像和个人运动数据，虽然当前项目主要在本地环境下演示，但若未来扩展到线上平台，就必须考虑视频存储周期、用户授权、数据脱敏、访问控制和日志审计等问题。特别是关键帧和模型原始响应中可能包含用户个人信息，系统需要提供删除机制和权限隔离机制，避免训练数据被无关角色查看。

最后，如果继续深入研究，还可以把规则评估与学习型方法结合起来。当前规则库依赖人工设定阈值，解释性较强但适应性有限；学习型模型可以从更多样本中学习动作质量差异，但解释性和数据需求更高。未来可以采用规则作为基础约束，再引入样本驱动的权重调整或个体化阈值学习，使系统既保持可解释性，又具备更强的适应能力。这样一来，本文当前完成的工程原型就可以进一步发展为更系统的运动姿态智能评估平台。

从应用推广角度看，系统后续还可以围绕“低门槛使用”继续优化。居家训练用户往往不会专门布置标准拍摄环境，因此系统需要在上传前给出拍摄提示，例如建议用户保持全身入镜、避免背光、尽量固定机位、选择动作持续最稳定的片段。这样的提示看似属于前端细节，实际上会直接影响姿态估计质量和最终反馈可信度。若输入视频质量更稳定，后端规则和模型反馈的有效性也会随之提高。

此外，未来也可以在报告生成上进一步完善。当前系统主要面向页面查看，后续可以把每次评估的关键帧、角度数据、问题列表和改进建议整理成可导出的训练报告，便于用户长期保存或与教练交流。报告中可以保留系统判断依据，同时避免过度堆叠技术细节，让普通用户读得懂、教练也能追溯分析依据。这样一来，系统就不只是一次性评估工具，而可以成为训练记录整理和阶段性复盘的辅助平台。

总之，本文已经完成了从问题提出、需求分析、系统设计、工程实现到测试验证的主要环节。后续工作可以继续沿着体式库扩展、样本积累、可视化优化、任务调度增强和隐私保护完善等方向推进。只要继续坚持“可运行、可解释、可追溯”的设计原则，当前原型仍有进一步扩展和深化的空间。

% 参考文献
\chapter*{参考文献}
\addcontentsline{toc}{chapter}{参考文献}

% 参考文献内容，不缩进，设置成字体：宋体，字号：五号，多倍行距1.25
{\songti\zihao{5}\setstretch{1.25}
\noindent
\sloppy
\begin{chineseenum}[label=\arabic*.]
  \item Lugaresi C, Tang J, Nash H, et al. MediaPipe: A Framework for Building Perception Pipelines[J]. arXiv preprint arXiv:1906.08172, 2019.
  \item Cao Z, Hidalgo G, Simon T, et al. OpenPose: Realtime Multi-Person 2D Pose Estimation Using Part Affinity Fields[J]. IEEE Transactions on Pattern Analysis and Machine Intelligence, 2021.
  \item Bradski G. The OpenCV Library[J]. Dr. Dobb's Journal of Software Tools, 2000.
  \item He K, Zhang X, Ren S, et al. Deep Residual Learning for Image Recognition[C]//Proceedings of the IEEE Conference on Computer Vision and Pattern Recognition. 2016: 770-778.
  \item Dosovitskiy A, Beyer L, Kolesnikov A, et al. An Image is Worth 16x16 Words: Transformers for Image Recognition at Scale[J]. arXiv preprint arXiv:2010.11929, 2020.
  \item Carreira J, Zisserman A. Quo Vadis, Action Recognition? A New Model and the Kinetics Dataset[C]//Proceedings of the IEEE Conference on Computer Vision and Pattern Recognition. 2017: 6299-6308.
  \item Redmon J, Farhadi A. YOLOv3: An Incremental Improvement[J]. arXiv preprint arXiv:1804.02767, 2018.
  \item Howard A G, Zhu M, Chen B, et al. MobileNets: Efficient Convolutional Neural Networks for Mobile Vision Applications[J]. arXiv preprint arXiv:1704.04861, 2017.
  \item Flask Documentation[EB/OL]. \url{https://flask.palletsprojects.com/}
  \item SQLite Documentation[EB/OL]. \url{https://www.sqlite.org/docs.html}
  \item MediaPipe Documentation[EB/OL]. \url{https://developers.google.com/mediapipe}
  \item RFC 7519. JSON Web Token (JWT)[EB/OL]. \url{https://datatracker.ietf.org/doc/html/rfc7519}
  \item Bishop C M. Pattern Recognition and Machine Learning[M]. Springer, 2006.
  \item Krizhevsky A, Sutskever I, Hinton G E. ImageNet Classification with Deep Convolutional Neural Networks[C]//Advances in Neural Information Processing Systems. 2012: 1097-1105.
  \item Chen T, Kornblith S, Norouzi M, et al. A Simple Framework for Contrastive Learning of Visual Representations[C]//Proceedings of the 37th International Conference on Machine Learning. 2020.
  \item Devlin J, Chang M W, Lee K, et al. BERT: Pre-training of Deep Bidirectional Transformers for Language Understanding[J]. arXiv preprint arXiv:1810.04805, 2018.
  \item Radford A, Narasimhan K, Salimans T, et al. Improving Language Understanding by Generative Pre-Training[EB/OL]. OpenAI, 2018.
  \item Brown T B, Mann B, Ryder N, et al. Language Models are Few-Shot Learners[J]. arXiv preprint arXiv:2005.14165, 2020.
  \item Bai J, Bai S, Chu Y, et al. Qwen Technical Report[EB/OL]. arXiv preprint, 2023.
  \item 阿里云百炼平台.千问多模态模型官方文档[EB/OL].
\end{chineseenum}
}

% 致谢
\chapter*{致谢}
\addcontentsline{toc}{chapter}{致谢}

% 正文各段首行缩进两个汉字，字体采用宋体小四号，行距为 1.25 倍
{\songti\zihao{-4}\setstretch{1.25}
\setlength{\parindent}{2em}
\indent 行文至此，几个月的毕设开发和论文写作终于要画上句号了。回顾整个过程，这不仅是一篇毕业论文的完成，更是一次从需求梳理、架构设计、编码实现到联调验证的完整工程实践，对于我个人的技术积累、问题分析能力和团队协作意识都具有重要意义。

\indent 首先，我想向我的指导老师张楗伟老师表示最诚挚的感谢。张老师从课题选择、研究方案、系统设计到论文修改的每一个阶段都给予了耐心指导。每当我在架构决策上犹豫不定、在技术实现上遇到瓶颈时，张老师总能及时指出关键点并提供可行建议，让我不断调整思路、逐步明确目标。尤其在论文写作的后期阶段，张老师多次对文稿提出专业性修改意见，帮助我把实验过程与工程实现更加清晰地呈现出来。

\indent 我还要感谢学院和实验室提供的良好学习环境与资源支持。无论是实验室的计算机设备、视频处理测试平台，还是论文写作期间的文献检索渠道，这些支持都为本课题的顺利推进提供了重要保障。感谢所有给予我协助的老师和工作人员，你们的帮助让我在本科学习阶段有机会完成这样一个相对完整的项目。

\indent 同时，我要感谢我的家人对我整个毕业设计期间的理解与支持。父母在生活和精神上的鼓励让我能够专注于项目开发，不必为日常事务过度分心。你们的陪伴和支持是我坚持下去的重要动力。

\indent 还要感谢我的同学和好友们，在项目开发与调试过程中给予我的宝贵意见。特别是同班同学在测试系统功能、提供使用反馈以及一起讨论算法方案时表现出的热情，极大提升了项目的实践价值。你们对问题的及时反馈和建议，使得系统在用户体验、异常处理和界面交互方面都有了更为稳妥的落地。

\indent 此外，我非常感谢开源社区与基础技术平台的支持。MediaPipe 提供了稳定的姿态估计能力，OpenCV 为视频处理和图像分析提供了可靠基础，Flask 为系统接口与服务调度提供了轻量级框架，SQLite 则使得本地持久化存储简洁可行。与此同时，千问多模态模型以及相关文档为本系统的反馈环节提供了重要参考，使得项目不仅停留在规则评估层面，还能向用户呈现更具可读性的建议文本。

\indent 最后，我要感谢所有曾经参与系统测试、提供问题报告和改进建议的志愿者。你们在不同场景、不同拍摄条件下的试用结果，为我发现了实际应用中容易出现的问题，比如视频上传稳定性、长轮询体验以及关键点可视化效果。正是你们的反馈，让我不断修正细节，使系统更贴近真实使用场景。

\indent 论文能够完成，离不开上述每一位给予我支持和关心的人。再次感谢所有在我本科毕业设计阶段给予帮助的老师、同学、家人和开源社区。你们的鼓励与支持将成为我继续前行的宝贵财富。
}
